\chapter{绪论}
\label{chap:introduction}

\section{研究背景与意义}
\subsection{多物理场有限元仿真的计算需求}
本文以多物理场有限元仿真为背景，围绕有限元系统构建与场间变量传递两个环节开展研究。
前者关注单元贡献向全局矩阵和载荷向量装配时的并发写入问题，
后者关注不同物理场采用非匹配网格时的传递误差与计算开销。
% TODO：补充工程应用背景及正式文献，说明两类问题对整体仿真流程的影响。

\subsection{并行装配的效率与正确性问题}
% TODO：说明共享自由度引起的写冲突、任务依赖与并行度之间的关系。
% TODO：用性能分析数据说明装配开销；不将局部装配加速等同于整体求解加速。
\thesisdraftnote

\subsection{非匹配网格的场变量传递问题}
% TODO：说明独立划分热网格与结构网格的需求，以及映射误差向结构响应的传播。
\thesisdraftnote

\section{研究问题与研究范围}
本文分别研究如何组织具有写入依赖的单元装配任务，以及如何根据局部场特征和网格尺度选择映射方法。
场映射接口面向节点型标量和多分量场；本文的实现与实验范围限定于线性四面体网格，
物理验证采用稳态单向热--结构耦合。
% TODO：区分接口可复用性与已完成的物理验证；不将热--结构结果直接推广到其他耦合方程。

\section{主要研究内容}
\subsection{基于空间分割前驱图的有限元并行装配}
% TODO：说明既有方法、本文拟改进的环节和对应的验证问题。
% TODO：研究资料指出该技术路线已有先例；必须明确本人工作的新增部分及引用归属。
\thesisdraftnote

\subsection{面向非匹配网格的自适应节点场映射}
研究内容包括空间定位与缓存复用、局部误差指示、映射策略选择，
以及与热--结构耦合方程相适应的等效热载荷虚功一致性校正。
最近邻、形函数插值、局部投影和径向基函数方法作为基础方法与对照方法。
% TODO：以实验结果界定策略的适用条件，避免预设自适应方法优于全部基线。

\section{论文组织结构}
第\ref{chap:foundations}章介绍相关研究与理论基础；
第\ref{chap:assembly}章讨论并行装配方法；
第\ref{chap:mapping}章讨论自适应场映射方法；
第\ref{chap:implementation}章介绍 MPIS 中的数据组织与实现；
第\ref{chap:experiments}章给出实验设计与结果分析框架；
第\ref{chap:conclusion}章总结研究工作并讨论后续方向。
